\section{转写校正与术语标准化}

本期没有平台字幕，正文基于 faster-whisper large-v3 CUDA 转写。数学和 AI 术语密集，自动转写会产生系统性误差，例如把 Lean 写成“令”或“林”，把 AI for Math 写成近似音，把 Carina/Hong/Verve Coffee 等专名误听，把 theorem proving、proof assistant、formalization 等术语拆散。生成讲义时必须做术语校正，否则读者会误解技术主线。

\subsection{术语消化：常见转写误差修正表}

\noindent\begin{longtable}{p{0.24\textwidth}p{0.32\textwidth}p{0.35\textwidth}}
\toprule
标准术语 & 可能的转写误差 & 校正理由 \\
\midrule
Carina Hong / 洪乐潼 & 洪乐同、洪乐童、Carina 音近误写 & 人名与 Axiom 背景相关，必须统一。 \\
Axiom & 公理、A型、Axium & 公司名；中文可解释为“公理”，但正文应保留英文。 \\
Lean & 令、林、line & proof assistant 名称，关系到形式化数学主线。 \\
AI for Math & Al for Math、AI for Mark & 本期核心方向。 \\
Proof Assistant & 证明助手、proof system 混用 & 指 Lean/Coq 这类验证系统。 \\
Formalization & 形式化、格式化混淆 & 指把数学写入形式系统，而不是排版格式化。 \\
Ken Ono & 小野肯、小野 Ken & 数学家专名，关联 Axiom 组织信号。 \\
Bounded / Free Attention & 注意力边界、自由注意力误译 & 这里是数学思考模式隐喻，不是 Transformer attention 机制本身。 \\
\bottomrule
\end{longtable}

\begin{warningbox}{自动转写不能直接当技术文本}
访谈类材料尤其容易在专名、论文名、系统名上误转写。高标准笔记必须基于章节、上下文和外部公开资料做术语校正，不能逐字照搬自动字幕。
\end{warningbox}

\subsection{本章小结}

转写只是原料，不是讲义。AI for Math 这种术语密集内容尤其需要先做术语标准化，再写结构化笔记。

